\documentclass[conference]{IEEEtran}
\usepackage{amsmath}
\usepackage{amssymb}
\usepackage{booktabs}
\usepackage{graphicx}
\usepackage{hyperref}

\title{Wavefront-64 Native Attention on CDNA4:\\
A Register-Topology-Aware Redesign of FlashAttention}

\author{
\IEEEauthorblockN{Chandan Kumar}
\IEEEauthorblockA{
Independent Research\\
Email: cs84861@gmail.com
}
}

\begin{document}

\maketitle

\begin{abstract}
FlashAttention and FlashAttention-3 achieve state-of-the-art performance on NVIDIA GPUs by optimizing memory movement and instruction scheduling.
However, these kernels were engineered under warp-32 execution assumptions.

On AMD CDNA4 (gfx950), execution occurs in Wavefront-64 mode with a fixed 256 VGPR budget per SIMD.
At HEAD\_DIM $\ge 128$, persistent full-head accumulation causes deterministic occupancy collapse.

This paper derives a formal register model for Wavefront-64, validates collapse empirically via compiler output, and proposes a partitioned-head redesign that restores occupancy without sacrificing MFMA throughput.
\end{abstract}

\section{Introduction}

Attention kernels dominate modern transformer workloads \cite{vaswani2017attention}.
FlashAttention \cite{dao2022flashattention} and FlashAttention-3 \cite{dao2024flash3} optimize memory traffic and scheduling for NVIDIA architectures.

However, GPU performance is constrained not only by memory bandwidth and instruction throughput, but by register topology.
We demonstrate that on CDNA4 Wavefront-64, persistent accumulator design induces a structural occupancy collapse at HEAD\_DIM = 128.

\section{CDNA4 Execution Model}

\subsection{Hardware Parameters}

\begin{table}[h]
\centering
\begin{tabular}{lc}
\toprule
Wavefront Size & 64 \\
VGPR per SIMD & 256 \\
SGPR per SIMD & 102 \\
LDS per CU & 160 KB \\
\bottomrule
\end{tabular}
\caption{CDNA4 (gfx950) architectural parameters}
\end{table}

\subsection{Occupancy Equation}

Let $R$ denote VGPR usage per lane.
Total SIMD VGPR budget = 256.

\begin{equation}
W = \left\lfloor \frac{256}{R} \right\rfloor
\end{equation}

where $W$ is waves per SIMD.

Critical boundary:

\begin{equation}
R \le 128 \Rightarrow W \ge 2
\end{equation}

\begin{equation}
R > 128 \Rightarrow W = 1
\end{equation}

Two waves are required to hide MFMA latency.

\section{Register Growth Model}

Let:

\begin{itemize}
\item $D$ = HEAD\_DIM
\item $R_{base}$ = non-accumulator registers
\item $R_{acc}$ = accumulator registers
\end{itemize}

Persistent full-head accumulation implies:

\begin{equation}
R_{total} = R_{base} + D
\end{equation}

Empirically:

\begin{equation}
R_{base} \approx 20
\end{equation}

For $D = 128$:

\begin{equation}
R_{total} \approx 148
\end{equation}

Thus:

\begin{equation}
W = \left\lfloor \frac{256}{148} \right\rfloor = 1
\end{equation}

Occupancy collapse.

\section{Empirical Validation}

RGA compiler output confirms:

\begin{table}[h]
\centering
\begin{tabular}{ccc}
\toprule
Accumulators & Used VGPR & Waves \\
\midrule
1 & 16 & 16 \\
8 & 80 & 3 \\
16 & 136 & 1 \\
32 & 296 & 1 \\
\bottomrule
\end{tabular}
\caption{Measured VGPR growth on gfx950}
\end{table}

Collapse observed at $\approx 130$ VGPR.

\section{Partitioned-Head Redesign}

We introduce partition factor $P$:

\begin{equation}
R_{total} = R_{base} + \frac{D}{P}
\end{equation}

For $D=128$, $P=2$:

\begin{equation}
R_{total} \approx 20 + 64 = 84
\end{equation}

\begin{equation}
W = \left\lfloor \frac{256}{84} \right\rfloor = 3
\end{equation}

Occupancy restored.

\section{When FlashAttention-3 Remains Optimal}

FlashAttention-3 \cite{dao2024flash3} remains optimal when:

\begin{itemize}
\item HEAD\_DIM $\le 64$
\item Kernel is memory-bound
\item FP8 reduces accumulator footprint
\end{itemize}

\section{Conclusion}

FlashAttention-3 optimizes scheduling and memory reuse.
However, Wavefront-64 introduces a different register topology constraint.

At HEAD\_DIM $\ge 128$, occupancy collapse is structural.

A partitioned-head Wavefront-64-native redesign restores wave residency and improves MFMA utilization.

This work reframes attention optimization as a register-topology alignment problem rather than a memory-only problem.

\bibliographystyle{IEEEtran}
\begin{thebibliography}{1}

\bibitem{vaswani2017attention}
A. Vaswani et al., ``Attention Is All You Need,'' 2017.

\bibitem{dao2022flashattention}
T. Dao et al., ``FlashAttention,'' 2022.

\bibitem{dao2024flash3}
T. Dao et al., ``FlashAttention-3,'' 2024.

\end{thebibliography}

\end{document}
